% sec:fd-convergence — Convergence Testing
\section{Convergence Testing}
\label{sec:fd-convergence}

A stencil derivation can be flawless and its truncation analysis
impeccable, yet the code that evaluates it may still contain an index
error, a misplaced sign, or a factor-of-two bug in the grid spacing.
Convergence testing catches these faults by turning the theoretical
predictions of \cref{sec:fd-truncation} into executable experiments: if
the code is correct, refining the grid must reduce the error at the rate
the theory predicts.  Once the grid is fine enough to be in the
asymptotic regime, any deviation is a reliable indicator that something
is wrong in the implementation --- making convergence tests the single
most powerful verification tool in numerical analysis \cite{Roache1998}.
\coderef{eh-compute}{stencils.rs}

\paragraph{Manufactured solutions.}

The simplest convergence test requires an exact answer to compare against.
For a single stencil operator, this is easy: choose a smooth function
$f(x)$ whose derivatives are known analytically, evaluate the stencil on
a grid of $N_x$ points, and measure the error at a representative
interior point.  The function $f(x) = \sin(2\pi x)$ on $[0, 1]$ serves
well: its derivatives are bounded, it exercises all stencil offsets
symmetrically, and its periodicity means that every grid point is
equivalent (no boundary contamination at the evaluation point).
\xrefnote{The worked examples of \cref{sec:fd-orders,sec:fd-truncation}
use this same test function, so the reader has already seen the errors
being measured here.}

The test proceeds as follows.  At each resolution $N_x$, compute the
stencil approximation $D_1[f]_i$ at the midpoint $i = N_x/2$ and compare
with the exact derivative $f'(x_i) = 2\pi\cos(2\pi x_i)$.  The absolute
error is
%
\begin{equation}
  \epsilon(N_x) = \abs{D_1[f]_i - f'(x_i)} \,.
  \label{eq:manufactured-error}
\end{equation}
%
For a $q$-th order stencil, the truncation error analysis of
\cref{sec:fd-truncation} predicts $\epsilon \propto \Delta x^q$, so
doubling the resolution (halving $\Delta x$) should reduce $\epsilon$ by
a factor of $2^q$.  The measured convergence ratio between successive
resolutions is
%
\begin{equation}
  \rho = \frac{\epsilon(N_x)}{\epsilon(2 N_x)} \,.
  \label{eq:convergence-ratio-test}
\end{equation}
%
Asymptotically, $\rho \to 2^q$ as $\Delta x \to 0$.  For fourth-order
stencils ($q = 4$), the prediction is $\rho = 16$.

\paragraph{A concrete test.}

The codebase implements exactly this procedure in the test
\texttt{test\_d1x\_sin\_convergence}, which runs at resolutions
$N_x = 16$, $32$, and $64$.  The test
asserts that the convergence ratio lies in the interval $[12, 20]$,
bracketing the theoretical value of $16$ with enough margin to absorb the
$1 + \order{\Delta x^2}$ relative correction from sub-leading truncation
terms (\cref{eq:convergence-ratio}).
\implnote{The acceptance window $[12, 20]$ is deliberately wide: the
test's purpose is to catch order-of-magnitude bugs (wrong stencil
coefficients, missing grid-spacing factors), not to measure the
truncation error coefficient to high precision.}

Running the test at the three resolutions gives errors of approximately
$2.0 \times 10^{-3}$, $1.3 \times 10^{-4}$, and $8.0 \times 10^{-6}$,
yielding ratios $\rho_{16 \to 32} \approx 15.7$ and
$\rho_{32 \to 64} \approx 15.9$ --- both within a few percent of $16$.
The agreement tightens as $N_x$ increases because the higher-order
correction terms shrink: the ratio approaches $2^q$ as
$\Delta x \to 0$.

\paragraph{Polynomial exactness tests.}

A complementary verification strategy exploits the polynomial exactness
of finite difference stencils.  The truncation error analysis of
\cref{sec:fd-truncation} shows that the leading error term involves
$f^{(p+q)}$, so a $q$-th order stencil for the $p$-th derivative
reproduces polynomials of degree $p + q - 1$ exactly --- every derivative
appearing in the error series vanishes for such polynomials.  Testing
this property is straightforward: evaluate the stencil on a polynomial
of the appropriate degree and check that the result matches the exact
derivative to machine precision.

For the fourth-order second derivative ($p = 2$, $q = 4$), any polynomial
of degree five or less has zero truncation error, so the stencil must
return the exact second derivative to machine precision.  The test
\texttt{test\_d2x\_quadratic\_exact} verifies this with $f(x) = x^2$:
the stencil returns $2.0$ to within $10^{-10}$, confirming that the
coefficients and grid-spacing factors are correct.
\warnnote{A polynomial exactness test passing does not guarantee
convergence order.  The stencil could have the correct coefficients but a
wrong power of $\Delta x$ in the denominator --- the polynomial test
wouldn't catch this because the exact result is independent of $\Delta x$.
Always combine polynomial tests with convergence-ratio tests.}

The mixed-derivative test \texttt{test\_d2xy\_polynomial} applies the
same principle to the separable stencil: with $f(x,y) = x^2 y^2$, the
exact mixed derivative $\partial^2 f / \partial x\, \partial y = 4xy$
must be recovered to machine precision.  This test validates both the
separable two-pass algorithm (\cref{sec:fd-stencils}) and its
implementation.

\paragraph{Self-convergence.}

Manufactured-solution tests work for isolated stencil operators, but the
full BSSN evolution has no closed-form solution to compare against.  The
self-convergence test addresses this by using the numerical solution
itself as the reference.

The idea is to run the same problem at three resolutions with grid
spacings in the ratio $4 : 2 : 1$.  Label the solutions by their
relative spacing: $u_4$ at the coarsest spacing $\Delta x$, $u_2$ at
$\Delta x / 2$, and $u_1$ at the finest spacing $\Delta x / 4$.  If the
dominant error is $\order{\Delta x^q}$, then the differences between
successive resolutions satisfy
%
\begin{equation}
  \frac{\norm{u_4 - u_2}}{\norm{u_2 - u_1}} = 2^q \,,
  \label{eq:self-convergence-ratio}
\end{equation}
%
where each difference is evaluated at the grid points of the coarsest
resolution, which are a subset of all three grids.  No exact solution
appears in this formula --- the method is self-referential, hence the
name.
\physnote{Self-convergence confirms that the code converges at the
expected rate, but not that it converges to the \emph{correct} answer
--- that requires a separate validation step, such as comparison with
perturbation theory or an independent code.}

For the BSSN equations evolved with fourth-order spatial stencils and a
fourth-order Runge--Kutta time integrator, the self-convergence ratio
should be $16$ in the spatial-error-dominated regime.  If the time
integrator has a different order from the spatial scheme, the ratio
reveals which error source dominates: for example, with a third-order
time integrator and CFL scaling $\Delta t \propto \Delta x$, a ratio
closer to $2^3 = 8$ would indicate that time-stepping error controls the
convergence (the interaction between spatial and temporal accuracy is
developed in \cref{ch:ode-integration}).  A ratio that doesn't match
any power of $2$ usually signals a bug rather than competing error sources.

\paragraph{Norm choice.}

The convergence ratio (\cref{eq:self-convergence-ratio}) depends on the
norm used to measure the solution difference.  Different norms weight
errors differently --- an averaged norm smooths over the grid while a
pointwise maximum exposes the worst offender --- and the choice can
affect whether a localised problem is visible.  Three standard choices
arise in practice:
%
\begin{equation}
  \norm{e}_1 = \frac{1}{N}\sum_{i=1}^{N} \abs{e_i} \,, \qquad
  \norm{e}_2 = \sqrt{\frac{1}{N}\sum_{i=1}^{N} e_i^2} \,, \qquad
  \norm{e}_\infty = \max_i \abs{e_i} \,,
  \label{eq:norms}
\end{equation}
%
where $e_i$ is the pointwise error (or solution difference) at grid point
$i$.  These are grid-function norms, scaled by the number of points so
that they approximate the corresponding continuous $L^p$ norms and remain
bounded as $N \to \infty$.  The $L^1$ and $L^2$ norms give a global
measure of accuracy; the $L^\infty$ norm reports the worst single point
and is useful for catching boundary contamination or isolated stencil
errors.  For smooth solutions, all three norms typically yield the same
convergence order.  Near discontinuities or sharp features, the
$L^\infty$ norm may show a lower order because a single poorly resolved
point dominates the maximum.  The codebase uses the $L^2$ norm as the
default for convergence monitoring.

\paragraph{Pitfalls.}

Three common mistakes can produce misleading convergence results.

First, evaluating the error at a point where the leading truncation
coefficient $C_q$ vanishes.  The convergence ratio
(\cref{eq:convergence-ratio}) requires $f^{(p+q)}_i \neq 0$; at a zero
of that derivative, the leading error term drops out and the ratio is
controlled by the next term in the series, producing a spuriously high
apparent order.  The $\sin(2\pi x)$ test function avoids this at
$x = 0.5$, where $f^{(5)}(0.5) = -(2\pi)^5 \neq 0$.

Second, using resolutions too coarse for the asymptotic regime.
At very low $N_x$, the higher-order correction terms in the truncation
error are comparable to the leading term, so the measured ratio can
differ substantially from $2^q$.  The Richardson extrapolation diagnostic
of \cref{sec:fd-truncation} can detect this: if the extrapolant differs
significantly from the fine-grid value, the grid is not yet in the
asymptotic regime and the reported convergence rate may be misleading.

Third, contamination from boundary treatment.  If the evaluation point is
close enough to a boundary that ghost zones or one-sided stencils
influence the result, the convergence order may reflect the lower-order
boundary scheme rather than the interior stencil.  The safest practice is
to evaluate at the domain midpoint, far from all boundaries.
\xrefnote{Boundary treatment and ghost-zone filling are discussed in
\cref{sec:fd-stencils}; the interaction with adaptive mesh refinement
boundaries is deferred to \cref{ch:amr}.}

\paragraph{The verification chain.}

Convergence testing forms the first link in a verification chain that
runs through the rest of this book.  Each subsequent chapter introduces a
new numerical component --- time integration
(\cref{ch:ode-integration}), mesh refinement (\cref{ch:amr}), constraint
evolution (\cref{ch:bssn-ccz4}) --- and each comes with its own
convergence tests that build on the methodology established here.  The
Kreiss--Oliger dissipation of \cref{sec:fd-dissipation}, for instance,
must not degrade the measured convergence rate; the order-preservation
analysis predicts this, and self-convergence tests confirm it
empirically.  The pattern is always the same: predict the convergence
rate from the truncation analysis, measure it computationally, and treat
any discrepancy as a bug until proven otherwise.  This discipline is not
optional for a numerical relativity code: without convergence testing,
there is no way to distinguish a physically meaningful result from a
numerical artefact.
